% Generated from locale/tr/content/first-order-logic/syntax-and-semantics/intro-semantics.tex; do not hand-edit.


\olfileid[tr]{fol}{syn}{its}

\olsection{Giriş}

İfadelere anlam vermek semantiğin alanıdır. Semantiğin merkezî kavramı,
!!a{structure} içinde sağlamadır. !!^a{structure}, dilin yapı taşlarına
anlam verir: !!a{domain}, boş olmayan bir nesneler kümesidir.
Niceleyiciler bu evren üzerinde dolaşacak biçimde yorumlanır;
!!{constant}s evrendeki elemanlara, !!{function}s !!{domain} içinden
yine kendisine giden fonksiyonlara, !!{predicate}s ise !!{domain}
üzerindeki bağıntılara atanır. !!{domain} ile temel söz dağarcığına
yapılan atamalar birlikte !!a{structure} oluşturur. !!^{variable}s,
!!{formula}s içinde geçebilir ve semantik verebilmek için bunlara da
!!{element}s öğelerini !!{domain} içinden atamamız gerekir; buna
değişken ataması denir. Son olarak sağlama bağıntısı bütün bunları bir
araya getirir. !!^a{formula}, !!a{structure}~$\Struct M$ içinde
!!a{variable} ataması~$s$'ye göre sağlanabilir; bunu
$\Sat M{!A}[s]$ biçiminde yazarız. Bu bağıntı da~$!A$'nın yapısı
üzerinde tümevarımla tanımlanır. Örneğin $(!A\land !B)$ ifadesinin
sağlanmasını,~$!A$ ile~$!B$ ifadelerinin sağlanmasına (ya da
sağlanmamasına) göre tanımlamak için mantıksal bağlaçların doğruluk
çizelgeleri kullanılır. Ardından !!{variable} atamasının,~$!A$
!!{formula} ifadesi !!a{sentence}, yani serbest değişken içermeyen bir
formülse önemsiz olduğu ortaya çıkar. Böylece
!!{sentence}s ifadelerinin !!{structure}s içinde doğrudan sağlandığından
(ya da sağlanmadığından) söz edebiliriz.

!!{sentence}s için $\Sat{M}{!A}$ sağlama bağıntısına dayanarak temel
semantik geçerlilik, mantıksal gerektirme ve sağlanabilirlik
kavramlarını tanımlayabiliriz. !!^a{sentence}, her !!{structure} onu
sağlıyorsa geçerlidir: $\Entails !A$. Bir !!{sentence}s kümesi,
$\Gamma \Entails !A$, ancak ve ancak $\Gamma$ içindeki bütün
!!{sentence}s ifadelerini sağlayan her !!{structure},~$!A$'yı da
sağlıyorsa bu tümceyi mantıksal olarak gerektirir. Bir
!!{sentence}s kümesi de, bir !!{structure} içindeki bütün
!!{sentence}s ifadelerini aynı anda sağlıyorsa sağlanabilirdir.
!!{formula}s tümevarımlı olarak tanımlandığından ve sağlama da tek tek
!!{formula}s ele alındığında bunların yapısı üzerinde tümevarımla
tanımlandığından, semantiğimizin özelliklerini ispatlamak ve tanımlanan
semantik kavramları ilişkilendirmek için tümevarımı kullanabiliriz.

